\section{DeepScholar-Bench：研究综合的三重难题}

\subsection{为什么 Related Work 是理想的长任务}

\videofigure{figures/deepscholar-bench.jpg}{DeepScholar-Bench 从研究综合角度补充 METR 的任务时长和 GDPval 的经济价值。}{00:51:48--00:53:04}

\videofigure{figures/deepscholar-task.jpg}{Benchmark 要求根据论文主题生成 Related Work，需要发现、筛选、组织和引用相关研究。}{00:53:04--00:54:52}

Related Work 没有唯一答案，却有明确质量结构：应覆盖关键工作、按主题组织、准确描述贡献、避免遗漏或伪造，并让引用可核验。这使它成为深度研究 agent 的高价值评估场景。

\begin{knowledgebox}{Research synthesis 不是普通长文本生成}
它包含检索 recall、来源判断、claim-to-source 对齐、跨文档综合和叙事组织。任一环节失败，最后段落都可能“读起来对，实际上不完整或不可证”。
\end{knowledgebox}

\subsection{三维评价}

\videofigure{figures/three-dimensional-eval.jpg}{评估同时覆盖知识综合、检索质量和引用可验证性。}{00:54:52--00:56:18}

三维指标分别回答：

\begin{itemize}
    \item \textbf{Knowledge synthesis}：结构是否连贯，事实是否正确，关键概念是否被组织成有意义的比较；
    \item \textbf{Retrieval quality}：来源是否相关、全面、重要，而非只找到若干容易搜索的论文；
    \item \textbf{Verifiability}：引用是否存在，引用是否支持对应 claim，出处是否可追溯。
\end{itemize}

\videofigure{figures/deepscholar-results.jpg}{不同系统在单项指标上各有优势，但没有系统能在全部指标上超过约 19\%。}{00:56:18--00:57:26}

\begin{importantbox}{最弱环节决定研究报告的可信度}
组织得好但来源不全，会形成漂亮的偏见；来源很多但 claim 对不上引用，会形成不可核验的堆砌。必须联合优化三维指标。
\end{importantbox}

\subsection{失败分析}

\videofigure{figures/deepscholar-failures.jpg}{主要失败来自来源覆盖不足、重要性判断低效、综合质量不稳和引用事实缺失。}{00:57:26--01:02:10}

模型往往能找到“相关”论文，却难以找到\textbf{全面且关键}的论文；也可能花大量 token 重复浏览边缘来源。即使检索到正确文献，系统还要识别每篇论文的贡献、限制和关系，最后把 claim 精确绑定到证据。

\begin{warningbox}{Citation precision 高不代表 coverage 高}
每条引用都真实且支持 claim，只说明没有明显伪造；如果漏掉半个领域的关键工作，综述仍然不合格。需要同时报告 precision、recall 与重要性加权覆盖。
\end{warningbox}

\begin{knowledgebox}{搜索停止条件是开放问题}
研究 agent 不知道“还有哪些重要论文没找到”。可用主题覆盖图、引用网络增益、来源重复率和新事实增益共同决定何时停止，而不是固定搜索十次。
\end{knowledgebox}

\subsection{本章小结}

DeepScholar-Bench 显示，长时程研究 agent 的主要瓶颈不是写一段流畅文字，而是高 recall 检索、重要性判断、跨文档综合和 claim-level verification。

